\section{Objections as denial conditions}
\label{sec:objections}

At this point the relevant objections are no longer free-floating. Each survives only by denying a
labeled definition, lemma, or theorem. Table~\ref{tab:strict-denial-table} states the denial burden in
its strict form before the representative objections are spelled out in brief.

\begin{longtable}{>{\raggedright\arraybackslash}p{0.28\textwidth} >{\raggedright\arraybackslash}p{0.44\textwidth} >{\raggedright\arraybackslash}p{0.18\textwidth}}
\caption{Strict denial table for the main hostile-referee objections.}\\
\toprule
\textbf{Objection} & \textbf{To maintain it, one must deny} & \textbf{Main location} \\
\midrule
\endfirsthead
\toprule
\textbf{Objection} & \textbf{To maintain it, one must deny} & \textbf{Main location} \\
\midrule
\endhead
The framework is optional & Theorem~\ref{thm:compressed-substrate-necessity} or the standing conditional fixed in Section~\ref{sec:targetclass}. & Sec.~\ref{sec:targetclass} \\
\addlinespace
You chose the quotient & Theorem~\ref{thm:equivalence-collapse}. & Sec.~\ref{sec:semanticcore} \\
\addlinespace
Alternative same-scope equivalence relation remains available & Theorem~\ref{thm:equivalence-collapse} together with Corollaries~\ref{cor:no-coarser-equivalence} and \ref{cor:no-finer-equivalence}. & Sec.~\ref{sec:semanticcore} \\
\addlinespace
Branch-local determinacy is enough & Theorem~\ref{thm:branchlocal-not-diachronic}. & Sec.~\ref{sec:branching} \\
\addlinespace
Everett still contains an implicit selector & Theorem~\ref{thm:selector-exhaustion} or Corollary~\ref{cor:bare-everett-exhaustion}. & Sec.~\ref{sec:taxonomy} \\
\addlinespace
Confirmation can be relativized or weakened & Theorems~\ref{thm:theory-level-confirmation-criterion} and \ref{thm:fixation-commutation-necessity}. & Sec.~\ref{sec:confirmation} \\
\addlinespace
Agency, memory, and anticipation are separate issues & Theorems~\ref{thm:diachronic-closure}, \ref{thm:identity-preserving-lineage}, and \ref{thm:diachronic-inheritance}. & Sec.~\ref{sec:diachronic_practice} \\
\addlinespace
There may be a fifth interpretive family & Theorems~\ref{thm:selector-locus-trichotomy} and \ref{thm:selector-exhaustion}. & Sec.~\ref{sec:taxonomy} \\
\addlinespace
An alternative selector outside the taxonomy remains available & Theorems~\ref{thm:selector-locus-trichotomy} and \ref{thm:selector-exhaustion}. & Sec.~\ref{sec:taxonomy} \\
\addlinespace
QBism or instrumentalism refutes the theorem ladder & Proposition~\ref{prop:deflation-exit}. & Sec.~\ref{sec:objections} \\
\bottomrule
\label{tab:strict-denial-table}
\end{longtable}

No row in Table~\ref{tab:strict-denial-table} is rhetorical. Each identifies the smallest formal denial
required to keep the corresponding objection alive on the fixed scope. The brief subsections that follow
spell out several representative rows.

\subsection{Objection 1: branch-relative determinacy is enough}

\begin{proposition}[Branch-relative determinacy does not suffice for the target discourse]
\label{prop:branch-relative-not-enough}
To maintain that branch-relative determinacy is sufficient, one must deny
Lemma~\ref{thm:branchlocal-not-diachronic} or restrict the target discourse so that it no longer asks
for theory-level diachronic reference.
\end{proposition}

\begin{proof}
Lemma~\ref{thm:branchlocal-not-diachronic} proves that branch-relative determinacy and diachronic
referential determinacy are distinct. The target explanatory practice of
Definition~\ref{def:target-practice} concerns the latter, not merely the former. So branch-relative
determinacy is enough only if one either denies that distinction or weakens the target discourse.
\end{proof}

\subsection{Objection 2: this is only self-location in different language}

\begin{proposition}[Self-location presupposes typed successor reference]
\label{prop:self-location-downstream}
To maintain that the problem is only self-location, one must deny
Theorem~\ref{thm:decision-theoretic-dependence} or show that self-locating discourse does not
presuppose typed successor and record reference.
\end{proposition}

\begin{proof}
Theorem~\ref{thm:decision-theoretic-dependence} shows that self-locating and
decision-theoretic discourse presuppose typed successor reference. So the self-location response is
not a prior solution unless that dependence is denied.
\end{proof}

\subsection{Objection 3: decision theory solves the problem}

\begin{proposition}[Decision theory presupposes rather than creates reference fixation]
\label{prop:decision-downstream-objection}
To maintain that decision theory solves the reference problem, one must deny
Theorem~\ref{thm:decision-theoretic-dependence} or show that a decision rule can itself create the
referential lineage to which it is applied.
\end{proposition}

\begin{proof}
A decision rule distributes weights or preferences only over already typed future alternatives.
Theorem~\ref{thm:decision-theoretic-dependence} formalizes exactly that dependence. So a
solution by decision theory alone requires denial of that theorem.
\end{proof}

\subsection{Objection 4: reference indeterminacy need not break science}

\begin{theorem}[Unresolved admissible divergence destabilizes theory-level science]
\label{thm:indeterminacy-breaks-science}
Under the standing conditional of Section~\ref{sec:targetclass}, unresolved admissible divergence in
reference assignments for measurement records makes theory-level confirmation, and with it the
memory/anticipation/agency package built around those records, assignment-relative rather than
theory-relative.
\end{theorem}

\begin{proof}
Corollary~\ref{thm:confirmational-consequence} proves that quotient-relevant divergence among
admissible assignments makes comparative confirmation assignment-relative. Theorem~\ref{thm:diachronic-practice}
proves that memory, anticipation, and agency inherit the same lineage dependence. So unresolved
admissible divergence does not leave the target scientific practice intact. It makes the central
uses of measurement records assignment-relative.
\end{proof}

\subsection{Objection 5: semantic infrastructure is unnecessary}

\begin{proposition}[Semantic infrastructure is required by the standing conditional]
\label{prop:not-adhoc}
Within the target class fixed by the standing conditional, a semantic infrastructure constraint is not an
optional embellishment. To deny it, one must deny
Theorem~\ref{thm:reference-stabilization-necessity} or reject the target explanatory practice.
\end{proposition}

\begin{proof}
Theorem~\ref{thm:reference-stabilization-necessity} states that a discourse-committing,
physically non-selective theory must restrict admissible reference assignments strongly enough for
fixation commutation to hold. A semantic infrastructure constraint is exactly such a restriction in the
non-dynamic, non-ontic case. So denying the need for semantic infrastructure requires denying the
theorem or exiting the target discourse.
\end{proof}

\subsection{Objection 6: the comparison with collapse and hidden-variable views is unfair}

\begin{proposition}[The comparative matrix is classificatory rather than adversarial]
\label{prop:fair-taxonomy}
The selector taxonomy does not privilege one interpretation family by rhetoric. To reject the
comparative matrix, one must identify a fifth same-scope locus of singular-reference fixing not covered
by Theorem~\ref{thm:selector-taxonomy}.
\end{proposition}

\begin{proof}
Theorem~\ref{thm:selector-taxonomy} classifies every package by the same structural question: where
is singular reference fixed, if it is fixed at all? Propositions~\ref{prop:collapse-dynamic},
\ref{prop:hidden-ontic}, \ref{prop:everett-semantic}, and \ref{prop:qbism-deflation} merely place
familiar interpretation families into that common matrix. So the classification is unfair only if a fifth
same-scope locus exists. The theorem denies exactly that.
\end{proof}

\subsection{Objection 7: the paper confuses semantics with physics}

\begin{proposition}[The result is conditional on explanatory use]
\label{prop:not-confusing-semantics-and-physics}
The paper does not replace physical analysis with semantic legislation. To maintain the opposite, one
must deny the conditional fixed in Section~\ref{sec:targetclass}: namely, that the theory is being used
to underwrite truth-apt diachronic discourse through measurement records.
\end{proposition}

\begin{proof}
Definitions and theorems in Sections~\ref{sec:targetclass}--\ref{sec:diachronic_practice} all take a
physical package as given and then ask what follows if that package is assigned a particular
explanatory burden. If that burden is removed, the conditional result does not apply. That is not a
confusion between semantics and physics; it is the explicit scope condition of the paper.
\end{proof}

\subsection{Objection 8: instrumentalism or QBism dissolves the problem}

\begin{proposition}[Discourse deflation exits the standing conditional]
\label{prop:deflation-exit}
Instrumentalist or QBist strategies do not refute the theorem ladder. They exit it by refusing the
standing conditional under which the paper runs.
\end{proposition}

\begin{proof}
Proposition~\ref{prop:qbism-deflation} places QBist and related instrumentalist strategies in class~(4)
of Theorem~\ref{thm:selector-taxonomy}: discourse deflation. Such strategies do not claim that the
theory underwrites the target class of truth-apt diachronic measurement discourse. Because the
standing conditional is refused, the need for reference stabilization inside the theory is likewise
refused. That makes the strategy a coherent exit from the target problem, not a counterexample to the
conditional theorem ladder.
\end{proof}

\subsection{What remains open}

Four things remain open and are consistent with the theorem ladder.
\begin{enumerate}
    \item The paper does not prove that a successful Everettian semantic-fixation strategy exists. It
    proves only that some such restriction is required if the target discourse is to remain
    theory-relative.
    \item The paper does not rank the interpretation families by independent metaphysical virtues.
    The taxonomy is structural.
    \item The paper does not deny that local branch-internal discourse may remain fully determinate
    once a branch is held fixed.
    \item The paper does not deny the coherence of discourse-deflationary responses. It says only that
    they are exits from the target explanatory burden rather than satisfactions of it.
\end{enumerate}
